\chapter{Zadatak 3 --- svođenje na standardni linearni program}
\label{ch:z3-lp}

\section{Što se ovdje radi}

Uklanja se apsolutna vrijednost iz funkcije cilja. To je jedini korak koji
problem od ,,konveksnog ali neglatkog`` čini linearnim programom.

\begin{predznanje}
\begin{itemize}[nosep,leftmargin=*]
  \item Diskretizirane veličine $h_i$, $D_1$, $D_2$ --- Poglavlje~\ref{ch:z3-diskretizacija}.
  \item Standardni oblik LP-a --- \S\ref{sec:z3-klasa}.
  \item Norma $\norm{\cdot}_1$ --- \S\ref{sec:vektori-matrice}.
\end{itemize}
\end{predznanje}

\section{Diskretni kriterij}
\label{sec:z3-kriterij}

\begin{korak}[Integral $\to$ suma]
Integral \eqref{eq:z3-kont} aproksimira se pravokutnom kvadraturom po čvorovima:
\begin{equation}
J_{\mathrm{disc}}(y) = \Delta x \sum_{i=0}^{N} \abs{y_i - h_i}
                     = \Delta x \, \norm{y - h}_1 .
\label{eq:z3-Jdisc}
\end{equation}
Faktor $\Delta x$ daje kriteriju dimenziju površine ($\mathrm{m}^2$). Kako je
riječ o \emph{pozitivnoj konstanti}, on ne mijenja minimizator --- samo mjerilo
kriterija. Bez njega bi $J$ za $N = 100$ bio pet puta veći nego za $N = 20$ pri
istom odstupanju, pa se vrijednosti ne bi mogle uspoređivati.
\end{korak}

\section{Trik s pomoćnim varijablama}
\label{sec:z3-l1-trik}

\begin{izvor}{Vježbe 4 --- slajd 5/40 (PDF str.~5)}
\slajd[0.86]{v4-norme-lp}

\textbf{Transkripcija.} \emph{Uobičajene norme:}
\begin{itemize}[leftmargin=*]
  \item $\norm{\cdot}_2$: Euklidska aproksimacija. Problem najmanjih kvadrata
        (\emph{least squares}). Rješenje $\xopt = (A\trans A)^{-1}A\trans b$ (uz
        pretpostavku da su stupci od $A$ linearno nezavisni).
  \item $\norm{\cdot}_\infty$: Chebysheva aproksimacija (minimax). Formulacija u
        obliku LP-a:
        \[
          \text{minimiziraj}_{x,t}\ t \quad \text{uz ogr.}\quad -t\mathbf{1} \leq Ax - b \leq t\mathbf{1},
        \]
        gdje je $t \in \R$ i $\mathbf{1}$ vektor stupac sa svim elementima
        jednakim $1$.
  \item $\norm{\cdot}_1$: suma apsolutnih vrijednosti reziduala. Formulacija u
        obliku LP-a:
        \[
          \text{minimiziraj}_{x,y}\ \mathbf{1}\trans y \quad \text{uz ogr.}\quad -y \leq Ax - b \leq y .
        \]
\end{itemize}

\textbf{Ovo je recept koji Zadatak 3 koristi doslovno.} Uoči razliku između
zadnja dva retka: kod $\ell_\infty$ pomoćna varijabla je \emph{jedan skalar} $t$
za sve reziduale (jer se minimizira najveći), a kod $\ell_1$ je \emph{vektor} ---
po jedna pomoćna varijabla za svaki rezidual (jer se zbrajaju svi).

\emph{Napomena o oznakama:} profesor pomoćni vektor zove $y$. U našem zadatku je
$y$ već zauzeto za visine nivelete, pa se pomoćni vektor ovdje zove $z$.
\end{izvor}

\begin{korak}[Zašto trik radi]
Uvodimo $z_i$ i tražimo $z_i \geq \abs{y_i - h_i}$, što je ekvivalentno paru
linearnih nejednakosti:
\begin{equation}
z_i \geq \abs{y_i - h_i}
\iff
\begin{cases}
\ \ y_i - h_i \leq z_i \\
-(y_i - h_i) \leq z_i
\end{cases}
\iff
\begin{cases}
\ \ y_i - z_i \leq h_i \\
-y_i - z_i \leq -h_i
\end{cases}
\label{eq:z3-par}
\end{equation}
Ključno je zapažanje: mi tražimo samo nejednakost $\geq$, a \emph{u optimumu
vrijedi jednakost}. Razlog: $z_i$ ulazi u funkciju cilja s pozitivnim
koeficijentom $\Delta x > 0$, pa je minimizacija tjera prema dolje, sve do donje
granice koju dopuštaju \eqref{eq:z3-par} --- a ta je granica upravo
$\abs{y_i - h_i}$. Zato je
\begin{equation}
\min \Delta x \sum_i z_i \quad\text{uz \eqref{eq:z3-par}}
\qquad\text{ima istu optimalnu vrijednost kao}\qquad
\min \Delta x \sum_i \abs{y_i - h_i}.
\end{equation}
\end{korak}

\begin{zamka}[Zamka: ,,zašto smijem $\geq$ umjesto $=$``]
Ovo je klasično potpitanje. Odgovor mora sadržavati riječ \emph{minimizacija}:
uvjet $z_i = \abs{y_i - h_i}$ nije linearan i ne bi se smio staviti u LP;
uvjet $z_i \geq \abs{y_i - h_i}$ jest linearan (kao par nejednakosti), a
jednakost se \emph{postiže sama} jer se $z_i$ minimizira. Da $z_i$ ulazi u cilj s
negativnim koeficijentom, ili da ne ulazi uopće, trik ne bi radio.
\end{zamka}

\begin{izvor}{Predavanje 5 --- LP i QP, slajd 24/34 (PDF str.~27)}
\slajd[0.86]{p05-po-dijelovima-afina}

\textbf{Transkripcija.} \emph{Primjer 1: Minimizacija po dijelovima afinih
funkcija.} Za zadan skup $\{a_i, b_i\}_{i=1,\dots,m}$ minimiziraj po dijelovima
afinu funkciju $f(x) = \max_{i} \bigl(a_i\trans x + b_i\bigr)$.
\emph{b) Formuliraj problem kao LP.} Ekvivalentan LP:
\[
\min_{x,t}\ t \quad \text{uz ogr.}\quad a_i\trans x + b_i \leq t, \quad i = 1,\dots,m .
\]

\textbf{Isti mehanizam, opći oblik.} Trik s $z_i$ je poseban slučaj ovoga:
$\abs{u} = \max\{u, -u\}$ je maksimum dviju afinih funkcija, pa se zamjenjuje
pomoćnom varijablom i parom nejednakosti. Ovaj se postupak zove
\emph{epigrafska reformulacija} --- umjesto same funkcije optimizira se preko
skupa iznad njezina grafa.
\end{izvor}

\section{Standardni oblik}
\label{sec:z3-standardni}

\begin{formula}{Linearni program Zadatka 3}
Vektor optimizacijskih varijabli je
\begin{equation}
\vartheta = \begin{bmatrix} y \\ z \end{bmatrix} \in \R^{2(N+1)},
\end{equation}
a linearni program glasi
\begin{equation}
\min_{\vartheta} \
\begin{bmatrix} \mathbf{0}\trans & \Delta x\,\mathbf{1}\trans \end{bmatrix} \vartheta
\quad \text{uz} \quad
A\vartheta \leq \beta, \quad A_{\mathrm{eq}}\vartheta = \beta_{\mathrm{eq}}, \quad z \geq 0,
\label{eq:z3-lp}
\end{equation}
gdje je
\begin{equation}
A = \begin{bmatrix}
I & -I \\ -I & -I \\ D_1 & 0 \\ -D_1 & 0 \\ D_2 & 0 \\ -D_2 & 0
\end{bmatrix},
\qquad
\beta = \begin{bmatrix}
h \\ -h \\ b_1\mathbf{1}_N \\ b_1\mathbf{1}_N \\ b_2\mathbf{1}_{N-1} \\ b_2\mathbf{1}_{N-1}
\end{bmatrix}.
\label{eq:z3-A}
\end{equation}
\end{formula}

\begin{korak}[Što je koji blok redaka]
\begin{enumerate}[leftmargin=*]
  \item $\begin{bmatrix} I & -I \end{bmatrix} \vartheta \leq h$ je
        $y_i - z_i \leq h_i$ --- prva polovica para \eqref{eq:z3-par}.
        $N+1$ redaka.
  \item $\begin{bmatrix} -I & -I \end{bmatrix} \vartheta \leq -h$ je
        $-y_i - z_i \leq -h_i$ --- druga polovica. $N+1$ redaka.
  \item $\pm D_1 y \leq b_1 \mathbf{1}$ je $\abs{\nabla y_i} \leq b_1$, rastavljeno
        na dvije nejednakosti. $2N$ redaka.
  \item $\pm D_2 y \leq b_2 \mathbf{1}$ je $\abs{\nabla^2 y_i} \leq b_2$.
        $2(N-1)$ redaka.
\end{enumerate}
Ukupno:
\begin{equation}
2(N+1) + 2N + 2(N-1) = 6N \ \text{nejednakosti}.
\end{equation}
Provjera: za $N = 20$ to je $120$, za $N = 100$ to je $600$ --- točno kako ispis
programa i navodi.
\end{korak}

\begin{korak}[Rubni uvjeti i granice varijabli]
\begin{itemize}[leftmargin=*]
  \item \textbf{Rubni uvjeti} $y_0 = a$ i $y_N = b$ su \emph{jednakosti}, pa idu
        u $A_{\mathrm{eq}}$, a ne među nejednakosti. Dva retka.
  \item \textbf{Varijable $y_i$ su slobodne} --- niveleta smije biti i ispod i
        iznad terena, pa donja granica mora biti $-\infty$. Ovo je česta greška:
        \texttt{linprog} po zadanome ne postavlja granice, ali ako se nesmotreno
        napiše \texttt{lb = zeros(...)}, niveleta se prisili da bude nenegativna
        i rješenje se tiho promijeni.
  \item \textbf{Varijable $z_i$ imaju donju granicu nula}, jer predstavljaju
        apsolutnu vrijednost. Formalno to slijedi već iz \eqref{eq:z3-par}
        (zbrajanjem dviju nejednakosti dobiva se $-2z_i \leq 0$), ali eksplicitna
        granica pomaže rješavaču.
\end{itemize}
\end{korak}

\begin{korak}[Veličina problema]
\begin{center}
\begin{tabular}{lcc}
\toprule
& $N = 20$ & $N = 100$ \\
\midrule
Čvorova $N+1$ & $21$ & $101$ \\
Varijabli $2(N+1)$ & $42$ & $202$ \\
Nejednakosti $6N$ & $120$ & $600$ \\
Jednakosti & $2$ & $2$ \\
\bottomrule
\end{tabular}
\end{center}
Oba su, po mjerilu s Predavanja 5 (,,LP i sa preko $10\,000$ varijabli``), vrlo
mali problemi.
\end{korak}

\begin{izgovor}
,,Apsolutnu vrijednost uklanjam po receptu s Vježbi 4: uvodim pomoćni vektor zet
i par nejednakosti, ipsilon i minus zet i manje ili jednako ha i, te minus
ipsilon i minus zet i manje ili jednako minus ha i. Traži se samo nejednakost, a
jednakost se u optimumu postiže sama, jer zet ulazi u funkciju cilja s
pozitivnim koeficijentom pa ga minimizacija tjera na donju granicu.

Vektor varijabli je ipsilon i zet zajedno, dakle dva puta en plus jedan
komponenti. Funkcija cilja je delta iks puta suma zetova --- ipsilon u nju uopće
ne ulazi.

Ograničenja: dva bloka za apsolutnu vrijednost, dva za nagib i dva za drugu
razliku, ukupno šest en nejednakosti. Rubni uvjeti idu kao dvije jednakosti.
Ipsilon je slobodna varijabla, zet ima donju granicu nula.``
\end{izgovor}
